% ---------------------------------------------------------------------------
% Payoff scaling shapes cooperation in LLM agents across languages
%
% Venue-neutral version of the Interface Focus manuscript.  The body text,
% figures, tables, appendix and bibliography are retained from the revision.
% Build the standalone documents with:
%   pdflatex main && bibtex main && pdflatex main && pdflatex main
%   pdflatex appendix && bibtex appendix && pdflatex appendix && pdflatex appendix
% ---------------------------------------------------------------------------
\documentclass[11pt,a4paper]{article}

\usepackage[T1]{fontenc}
\usepackage[utf8]{inputenc}
\usepackage{lmodern}
\usepackage[margin=1in]{geometry}
\usepackage{graphicx}
\usepackage{amsmath,amssymb}
\usepackage{booktabs}
\usepackage{tabularx}
\usepackage{caption}
\usepackage{authblk}
\usepackage[numbers,sort&compress]{natbib}
\usepackage{microtype}
\usepackage{float}
\usepackage{placeins}
\usepackage{enumerate}
\usepackage{xcolor}
\usepackage[colorlinks=true,allcolors=blue!55!black]{hyperref}
\usepackage{xr}
\externaldocument[][nocite]{appendix}

\captionsetup{font=small,labelfont=bf,width=\textwidth}
\setlength{\parskip}{0pt}
\renewcommand{\baselinestretch}{1.15}
\raggedbottom

% Keep substantial figures at the top of a text page instead of creating
% sparsely populated float pages in the neutral article layout.
\setcounter{topnumber}{2}
\setcounter{bottomnumber}{2}
\renewcommand{\topfraction}{0.95}
\renewcommand{\bottomfraction}{0.80}
\renewcommand{\textfraction}{0.05}
\renewcommand{\floatpagefraction}{0.80}

\newcommand{\lam}{\lambda}
\newcommand{\ci}[2]{(95\% CI #1, #2)}

% In the preamble
\usepackage{authblk}
\usepackage{hyperref}

\title{Payoff scaling shapes cooperation in LLM agents across languages}

\author[]{%
The Anh Han\textsuperscript{1\textdagger*},
Trung-Kiet Huynh\textsuperscript{2,4\textdagger},
Dao-Sy Duy-Minh\textsuperscript{2,4\textdagger},
Thanh-Bang Cao\textsuperscript{3,4},
Phong-Hao Le\textsuperscript{3,4},
Hong-Dan Nguyen\textsuperscript{3,4},
Phu-Quy Nguyen-Lam\textsuperscript{2,4},
Minh-Luan Nguyen-Vo\textsuperscript{3,4},
Hong-Phat Pham\textsuperscript{3,4},
Phu-Hoa Pham\textsuperscript{2,4},
Thien-Kim Than\textsuperscript{3,4},
Chi-Nguyen Tran\textsuperscript{3,4},
Huy Tran\textsuperscript{3,4},
Gia-Thoai Tran-Le\textsuperscript{3,4},
Alessio Buscemi\textsuperscript{5},
and Le Hong Trang\textsuperscript{3,4*}
}

\affil[]{%
\textsuperscript{1}School of Computing, Engineering and Digital Technologies,
Teesside University, Middlesbrough, United Kingdom\\
\textsuperscript{2}Faculty of Information Technology,
University of Science (HCMUS), Ho Chi Minh City, Vietnam\\
\textsuperscript{3}Faculty of Computer Science and Engineering,
Ho Chi Minh City University of Technology (HCMUT),
Ho Chi Minh City, Vietnam\\
\textsuperscript{4}Vietnam National University,
Ho Chi Minh City (VNU-HCM), Ho Chi Minh City, Vietnam\\
\textsuperscript{5}Luxembourg Institute of Science and Technology (LIST),
Luxembourg
}

\renewcommand\Authfont{\normalsize}
\renewcommand\Affilfont{\small\normalfont}
\setlength{\affilsep}{0.5em}

\date{%
\vspace{-1.0em}
\footnotesize
\textsuperscript{\textdagger}\,These authors contributed equally to this work.\\
\textsuperscript{*}\,Corresponding authors:
\href{mailto:t.han@tees.ac.uk}{t.han@tees.ac.uk}
and
\href{mailto:lhtrang@hcmut.edu.vn}{lhtrang@hcmut.edu.vn}.
}

\begin{document}

\maketitle

\begin{abstract}
Large language models (LLMs) increasingly act as autonomous agents negotiating and deciding for human users, making their cooperative dispositions a central question for AI governance. We examine how two fundamental levers shape cooperation: the nominal magnitude of the stakes and the language in which the interaction is framed. In classical game theory, multiplying all payoffs by a positive constant is an invariant transformation that leaves preference rankings, best replies, and equilibria untouched. We show, across five frontier LLMs playing an iterated prisoner's dilemma across ten payoff scales and five languages, this invariance is systematically violated. Rescaling stakes alters cooperation rates dramatically, driving divergent shifts across models that cancel out when averaged. Scale sensitivity also varies across languages. Payoff scale also changes persona compliance: at sub-unit stakes, all five models cooperate more under the selfish than the cooperative persona; in Claude and Gemini this effect reverses sign above the sub-unit regime. Across the evaluated models, own-action conditioning is stronger than partner-action conditioning. These findings show that cooperation metrics cannot be treated as intrinsic model traits without specifying payoff scale and language.
\end{abstract}

\noindent\textbf{Keywords:} machine behaviour; large language models; iterated prisoner's dilemma; payoff scaling; utility invariance; cross-lingual behaviour


\section{Introduction}\label{sec:intro}

Large language models (LLM) are increasingly deployed as agents that
negotiate, coordinate and decide on behalf of their users, in recommendation
systems, negotiation tools and multi-agent
assistants~\cite{park2023generative,tessler2024ai,hammond2025multi,cook2024llm}.
In these settings, behaviour emerges from strategic interaction rather than from
isolated model outputs, and the agents face cooperation dilemmas with
solutions  directly relevant to  safety, coordination and AI
governance~\cite{dafoe2021cooperative,lu2024llms,han2026social,idowu2026mapping}.
A fast-growing literature measures how they behave when they do, using methods from behavioural economics and cognitive
psychology~\cite{binz2023using,horton2023large,aher2023using,rilla2026recognising} and  the
classical games~\cite{mei2024turing}. Equally consequential is how far that
behaviour adapts to the costs and benefits at issue, since a deployed agent
meets the same dilemma at many different magnitudes~\cite{hammond2025multi}.

A growing body of work evaluates language models through game-theoretic
lenses, using evolutionary and repeated games, reporting
systematic departures from equilibrium outcomes, persistent cooperative biases and
sensitivity to contextual factors such as language and
incentives~\cite{buscemi2025fairgame,sun2025gtllm_survey,mao2025alympics,pal2026strategiescooperationdefectionlarge,willis2025will,fan2024can,pires2025large,fu2026optimal,xie2026evolutionary}.
The prisoner's dilemma has been a well adopted framework ~\cite{axelrod1981evolution,axelrod1984evolution,nowak1992tit,nowak1993strategy,nowak2006five,pressDyson2012}.
Our study is closest to FAIRGAME, which introduced a game-theoretic
evaluation framework and reported payoff and language sensitivity in repeated
games~\cite{buscemi2025fairgame}. We retain its no-communication prisoner's
dilemma protocol so that the comparison is anchored to the same task, while
adding a positive-rescaling null spanning ten scales, five languages and five
model endpoints. This makes our
contribution a test of invariance and heterogeneity, rather than an evaluation of LLM behaviours in a new strategic interaction.

Prior work showed that the LLM models differ
in how readily they retaliate~\cite{akata2025} or 
forgive~\cite{fontana2025nicer}, how much weight they give the game rather than
its framing  ~\cite{loreheydari2024strategic} and how far an instructed
disposition carries into conditioned
reciprocity~\cite{phelps2023investigating}. These studies, however, share two largely unremarked reporting conventions. First, they characterise a model using a cooperation rate, measured under a single payoff matrix. Second, they report what agents did, rather than the decision rule that generated the observed behaviour. This makes it difficult to tell whether a reported cooperation rate is a stable property of the model, an artefact of the particular payoff scale and framing, or the aggregate trace of a different strategic rule.

Evaluating strategic behaviour at the level required for alignment and governance therefore requires more than aggregate outcomes. Drawing on
behavioural and evolutionary game theory~\cite{sigmund:2010bo}, a behavioural
intention can be conceptualised as an agent's \emph{strategy}: a decision rule
mapping an interaction history to the next
action~\cite{han2013intention,diStefanoIntention2023,Han2012ALIFEjournal,fujimoto2019functional}.
The canonical strategies of repeated games, e.g. Always Cooperate, Always Defect,
Tit-for-Tat and Win-Stay-Lose-Shift~\cite{Axelrod1980Effective,sigmund:2010bo},
provide an interpretable vocabulary for such intentions, and prior work has
shown that they can be inferred from noisy behavioural trajectories by
supervised learning~\cite{han2013intention,diStefanoIntention2023} and extended by
probabilistic models to mixed or stochastic play~\cite{inferenceStrategies}.
The inference is not free. Separating the memory-one rules from observed play
has been shown to need a long run of observations even where the resident rule
is fixed by construction~\cite{kim2021win,hilbe2013memory}, which is why the
read-out used here is reported as a resemblance, with its provenance beside it
(Section \ref{sec:readout}).

In parallel, a language model's behaviour varies with the language in which the
prompt is written ~\cite{buscemi2025fairgame}. Strategic reasoning, risk
sensitivity and cooperation patterns depend on linguistic framing, sometimes by
margins comparable to the differences between
models~\cite{loreheydari2024strategic,buscemi2025llms,buscemi2025strategic}.
Yet it remains unclear whether the strategies that agents appear to play are
themselves language-dependent, and whether such dependence introduces systematic
bias into multi-agent settings.




This dependence is an instance of a more general problem: measurements of
language-model behaviour are often sensitive to aspects of the elicitation
procedure that are not themselves the intended object of measurement. Accuracy
can vary with the ordering of in-context examples~\cite{lu2022fantastically},
and with formatting choices that preserve the semantic content of the prompt,
by margins large enough to obscure the differences between the models under
comparison~\cite{sclar2024quantifying}. Persona instructions are similarly
consequential: they alter model outputs~\cite{deshpande2023toxicity}, affect
measured capability~\cite{salewski2023context}, and, when specified through a
demographic backstory, can induce response distributions resembling those of
the corresponding human subgroup~\cite{argyle2023out}. They therefore do not
establish a well-defined null condition against which the effect of prompting
can be assessed.


The choice of units (of measurement)  is rarely presented as a design decision, because
in the standard theory it is not one. A payoff matrix is a von Neumann-Morgenstern
utility, defined only up to a positive affine
transformation~\cite{von2007theory,luce1957games}. Multiplying every
entry in the payoff matrix by ten changes no preference ordering, no best response, no
equilibrium, and no trajectory of the replicator dynamics. The result is not a
new game with larger payoffs, but the same game written in different units.  The
theory of the iterated dilemma inherits this invariance: its predictions depend 
on the ordering of the  payoff-matrix entries  and on ratios among them, rather than on
their absolute magnitudes~\cite{axelrod1981evolution,nowak1992tit,nowak1993strategy,pressDyson2012}.
For this reason, the modelling literature often normalises the payoff scale
away~\cite{szabo2007evolutionary,perc2010coevolutionary,wang2015universal}.
Moreover, Recent theory strengthens this point: in models in which social
learners update on the payoff they most recently received, the condition for
conditional cooperation to persist under strong selection contains neither the
benefit nor the cost of cooperation, and depends only on the ordinal structure
of the payoffs~\cite{glynatsi2024evolution}.

The invariance is a property of the theory, not of human decision-making. Human valuation is neither
linear in magnitude nor anchored at zero~\cite{kahneman1979prospect}, human
choice reverses between logically equivalent descriptions of one
problem~\cite{tversky1981framing}, and human cooperation in the dilemma itself
responds to the size of the stakes~\cite{krockow2016cooperation,list2006friend}.
That gap has been argued to be structural rather than incidental, with the
linguistic description of an action entering the utility function directly and
a shift proposed from outcome-based to language-based
preferences~\cite{capraro2024language}.

Two things are established, and they have not been put together. The first is
that what a language model does depends on features of the prompt that carry no
task content: the order of the examples, the formatting, the persona, the
language. The second is that multiplying every payoff-matrix entry by a positive constant
carries no task content in the strictest sense available, because the matrix it
produces is  the same game. What is not known is whether
the strategic behaviour of a language model respects that invariance. To our
knowledge, no prior study has formulated positive payoff rescaling as a
theorem-backed zero-effect invariance test across a broad multi-model,
multilingual grid. Every manipulation in the sensitivity literature could, in principle, induce legitimate behavioural changes.  
Consequently, observed changes under these manipulations are difficult to interpret, since the expected direction or magnitude of the effect is typically unspecified. Such changes may therefore be treated either as methodological noise or as empirical findings, but only relative to an implicit and under-theorised baseline. Payoff rescaling is distinctive because its theoretical prediction is explicit: behaviour should remain unchanged. A measured departure from this prediction therefore has a clearly defined null reference, rather than only a contrast with another experimental condition.

% What follows if the invariance fails is a question about what a published
% cooperation rate means. Machine behaviour proposes that autonomous systems be
% studied empirically, with the instruments of the behavioural sciences, because
% what such a system does is not deducible from how it was
% built~\cite{rahwan2019machine}. Those instruments arrive with a known
% difficulty: a behavioural manipulation almost never carries a predicted effect,
% so a measured change has to be considered against a comparison condition the
% experimenter chose. The invariance classes of decision theory are the
% exception, supplying a manipulation whose predicted effect is a stated number,
% and the number is zero. The consequences run through the strands of machine
% behaviour this study touches. The persona line is the one place where the
% experimenter addresses the agent about its social role, and whether it takes
% hold turns out to depend on the units the matrix was shown in. Models of
% mixed human-machine populations carry the cooperativeness of the artificial agent as
% a parameter and ask where the population
% settles~\cite{santos2018social,han2011intention,han2013good,perc2017statistical,paiva2018engineering};
% a rate partly determined by the experimenter's units is not the quantity those
% models need.



Taken together, these considerations motivate three questions about the effects of payoff rescaling:
\begin{enumerate}[(1)]
\item Is strategic behaviour invariant to   rescaling of payoffs, and how large are any deviations from invariance?
\item Does the result depend on the language of the prompt and on the persona assigned to the model?

\item Does payoff rescaling affect only how often the model  cooperates, or does it also alter the underlying strategy it adopts?

\end{enumerate}
We answer them on a complete and exactly balanced grid of 10{,}000 dyads
(Section \ref{sec:methodology}), and compare the responses against an evolutionary game
baseline computed on the same grid.

%Here, we turn a theorem into an instrument. 
The main contributions of the present
paper are: \textbf{(i)} A payoff-rescaling protocol whose predicted behavioural
effect is exactly zero, demonstrating that five frontier language models
systematically violate scale invariance (Section \ref{sec:payoff_magnitude}), with
cooperation rates shifting substantially across stakes. \textbf{(ii)} A
per-model analysis revealing that scale departures diverge sharply across
models to the point of anticorrelation, so that pooled analyses produce
an artificial average that describes no real model: variance is dominated by
model individuality and model-by-scale interactions. In the five-model
cell-mean decomposition, the shared scale main effect is largely carried by
Grok, whereas the dyad-level pooled association persists in leave-one-model-out
fits. \textbf{(iii)} Cross-lingual and structural evidence showing that scale
sensitivity depends heavily on the prompt language, gates persona compliance
at sub-unit stakes, and alters conformity to canonical strategies.
\textbf{(iv)} An analytical evolutionary game baseline that moves in the
opposite direction to the pooled empirical corpus. The transition analysis
shows stronger persistence than reciprocity, without identifying a causal
mechanism for cooperation (Section \ref{sec:conditioning}, appendix~\ref{appendix:conditioning_ci}).

% %An earlier and smaller collection on this project has been reported separately
% by the present authors~\cite{selfpreprint2026scaling}: threemodels at three
% payoff scales, in the same five prompt languages, with cooperation read off the
% trajectories and reported as rising with the stakes. The present article
% differs in what it collects and in what it claims. It collects five models
% rather than three and ten payoff scales rather than three, spanning five orders
% of magnitude, on a grid whose completeness and balance are enforced by the
% analysis build rather than assumed. It frames the rescaling as an invariance
% whose predicted effect is a theorem rather than as a manipulation of the stakes
% whose effect was expected. And it reports every quantity per model, which is
% what shows the movement to be heterogeneous, with two of the five running
% downwards, and the pooled statement to be an artefact of averaging.


% ===========================================================================
\section{Material and methods}\label{sec:methodology}

\subsection{Framework}\label{sec:background}

FAIRGAME, the Framework for AI Agents Bias Recognition using Game
Theory~\cite{buscemi2025fairgame}, supplies the computational infrastructure
for this study. It supports systematic and reproducible evaluation of language
models through controlled game-theoretic experiments: an experimental condition
is declared in a configuration file that fixes the payoff structure, the
horizon, the model backend and the prompt language; at run time the framework
combines that declaration with language-specific prompt templates, simulates
the repeated normal-form game, and logs the round-by-round trajectory. We
extend it in three ways. First, with a payoff-scaling factor applied to the
prisoner's dilemma, which is the manipulation this paper examines 
(Section \ref{sec:gameImportance}). Second,  we apply a completeness safeguard that prevents the release of an analysis table unless the experimental grid is exactly balanced and the recorded payoffs for each game correctly recover the payoff scale assigned to its corresponding cell (i.e. a  combination of experimental conditions, such as model, language, persona, and payoff scale). Third, we include a strategy read-out that identifies which canonical memory-one rule each trajectory most closely resembles, while also reporting the provenance of that classification (Section \ref{sec:readout}).

The design is a complete factorial crossing of five language models, ten payoff
scales, five prompt languages, four persona pairings and ten replicates. The
unit of collection is the dyad, one ten-round game between two agents of the
same model, contributing two agent-games, an agent-game being the ten decisions
one agent takes in one game. Each of the 250 model-by-language-by-scale cells
holds exactly 40 dyads, the four persona pairings by ten replicates; pooling
the five languages gives 200 dyads and 400 agent-games in each of the 50
model-by-scale cells. The main-text corpus is 10{,}000 dyads, 20{,}000
agent-games and 200{,}000 decisions, and is the grid every main-text figure and
table is computed on. Balance is enforced rather than assumed
(appendix~\ref{appendix:corpus}).

\subsection{The payoff-scaled prisoner's dilemma and its invariance}\label{sec:gameImportance}

Both agents receive the prisoner's-dilemma prompt of the FAIRGAME
protocol~\cite{buscemi2025fairgame} in its no-communication, announced-horizon
condition. The scenario casts the strategic interaction as a criminal
interrogation, states the four payoff outcomes as penalties, names the two
available actions Option~A and Option~B, carries the running history of prior
moves and realised penalties, announces a fixed horizon of $T = 10$ rounds, and
instructs each agent to minimise its cumulative penalty.

The row player's baseline payoff matrix for the Prisoner's Dilemma is
\begin{center}
\small
\begin{tabular}{c|cc}
 & Option~A & Option~B \\
\hline
Option~A & $(6, 6)$ & $(0, 10)$ \\
Option~B & $(10, 0)$ & $(2, 2)$
\end{tabular}
\end{center}
where the first and second entries in each cell denote the penalty payoffs to
the row and column player, respectively. In standard game-theoretic notation,
this matrix assigns the four canonical outcomes:
\[
T = 0, \quad R = 2, \quad P = 6, \quad S = 10,
\]
representing the temptation ($T$), reward ($R$), punishment ($P$), and sucker
($S$) penalty payoffs. The penalty ordering
\[
T < R < P < S \quad \text{with} \quad 2R < T + S
\]
is the exact isomorphic image of the familiar canonical Prisoner's Dilemma
condition $T_{\text{util}} > R_{\text{util}} > P_{\text{util}} > S_{\text{util}}$
and $2R_{\text{util}} > T_{\text{util}} + S_{\text{util}}$ defined on utilities
$u = -\text{penalty}$:
\[
T_{\text{util}} = 0 > R_{\text{util}} = -2 > P_{\text{util}} = -6 > S_{\text{util}} = -10.
\]
Because lower penalties correspond to higher utilities, Option~A strictly
dominates Option~B in every column (facing Option~A, choosing A yields penalty 6
versus 10; facing Option~B, choosing A yields penalty 0 versus 2). Option~A is
therefore the defecting action, while mutual Option~B yields penalty 2 versus
penalty 6 for mutual Option~A, making Option~B the cooperative action;
throughout this paper, cooperation denotes choosing Option~B.

The realised score-action relationships reproduce the intended dilemma
structure in the logs: own cooperation correlates with own utility at Spearman
$\rho = -0.015$, while the opponent's cooperation correlates at
$\rho = +0.916$, and dyad welfare correlates with joint cooperation at
$\rho = +0.985$. These checks validate the payoff mapping and action polarity,
not model comprehension of the dilemma (appendix~\ref{appendix:corpus}).

To manipulate the stakes of the game without altering its strategic structure,
best-reply correspondences, or Nash equilibria, we introduce a positive scalar
parameter $\lam > 0$ and multiply all four baseline cells by $\lam$
(figure~\ref{fig:overview}b). In our experiments, we evaluate ten values spanning
five orders of magnitude:
\[
\lam \in \{0.01,\ 0.10,\ 0.25,\ 0.50,\ 1.0,\ 2.0,\ 5.0,\ 10.0,\ 100.0,\ 1000.0\},
\]
corresponding to attenuated ($\lam < 1.0$), baseline ($\lam = 1.0$), and
amplified ($\lam > 1.0$) payoff regimes. For example, at the sub-unit scale
$\lam = 0.10$, the row player's payoff matrix becomes
\begin{center}
\small
\begin{tabular}{c|cc}
 & Option~A & Option~B \\
\hline
Option~A & $(0.6, 0.6)$ & $(0, 1.0)$ \\
Option~B & $(1.0, 0)$ & $(0.2, 0.2)$
\end{tabular}
\end{center}
where every  payoff entry is at most 1.0. At the amplified scale $\lam = 10.0$,
it becomes
\begin{center}
\small
\begin{tabular}{c|cc}
 & Option~A & Option~B \\
\hline
Option~A & $(60, 60)$ & $(0, 100)$ \\
Option~B & $(100, 0)$ & $(20, 20)$
\end{tabular}
\end{center}
with the ordering $T < R < P < S$ strictly preserved in all cases. The complete
matrix for all ten scales is provided in appendix~\ref{appendix:payoff_scaling}.

Under classical decision theory and game theory, multiplying utilities by
$\lam > 0$ constitutes a positive affine transformation of a
von Neumann-Morgenstern utility function~\cite{von2007theory,luce1957games}:
\[
u'(a) = \lam u(a) + c \quad (\lam > 0,\ c = 0).
\]
Such transformations preserve all preference orderings over lotteries, best
replies, Nash equilibria, and deterministic replicator dynamics orbits.
Consequently, the normative theoretical prediction for the effect of $\lam$ on
rational agents is exactly zero. This construction isolates the effect of payoff
magnitude from structural incentives: any systematic departure from invariance
constitutes an empirical discovery about machine behaviour rather than a
rational adaptation to changing game-theoretic payoffs.

\begin{figure}[!htbp]
\centering
\includegraphics[width=\linewidth]{figures/prisoners_dilemma_overview.pdf}
\caption{\textbf{Overview of the repeated prisoner's-dilemma experiment.}
Two language-model agents from the same provider repeatedly play the stage game across ten payoff scales, with penalties to minimise. The main-text corpus covers five frontier models, five prompt languages, four persona pairings and 10{,}000 dyads (20{,}000 agent-games and 200{,}000 decisions).}
\label{fig:overview}
\end{figure}

\subsection{Models, prompt languages, personas and protocol}\label{sec:protocol}

Five frontier models were run, each against itself; two different models were
never paired. They were Claude Haiku 4.5, GPT-5.4 Nano, Gemini 3.5 Flash-Lite,
Qwen3 235B-A22B and Grok 4.20. Their model identifiers are listed in
appendix~\ref{appendix:endpoints}.

The prompt languages are English, French, Vietnamese, Chinese and Arabic, in
templates byte-identical to the FAIRGAME resource files and checked against
them before every collection run. Artefacts of those files are reproduced
rather than repaired, the consequential one being that the Arabic and Chinese
versions state the four cells as penalties but phrase the goal sentence as
maximising a reward, so we treat the five languages as five distinct stimuli
rather than five renderings of one (appendix~\ref{appendix:templates}). Each
agent also receives one line naming its persona, in the language of the prompt:
in English, \emph{You are cooperative.} or \emph{You are selfish.} The two are
crossed in all four ordered pairings, ten replicates each. The opponent's
persona is never disclosed, so a pairing reaches a player only through realised
behaviour and not before round 2, and we treat it as a design factor and a
regression covariate rather than as information the agent held.

Every game runs ten rounds with simultaneous moves: both agents see the same
history, rounds 1 to $r-1$, and the round is scored once both have replied.
Sampling uses documented matching by game identifier across payoff scales.
The released tables do not retain the API seed values, so this record verifies
the matching structure but not provider-specific seed semantics. Decisions are elicited
as free text and parsed by the protocol's own matcher. A reply it cannot
resolve is regenerated, up to two further attempts, and if all three fail the
decision is recorded as Option~A, which under this framing is defection: the
fallback can only push a measured cooperation rate down, never up, and it
cannot manufacture the cooperative behaviour reported here. Each collection run
asserts that its fallback rate is at most 0.02, but the achieved rates are not
carried in the released per-game tables, so we cannot report how they vary with
the payoff scale (appendix~\ref{appendix:endpoints}).

\subsection{Strategy read-out}\label{sec:readout}

While the framework provides complete gameplay trajectories and descriptive
metrics such as cooperation rates, these capture what agents do rather than the
decision rule behind it. Following prior work on inferring canonical strategies
from noisy repeated-game data~\cite{han2013intention,Han2012ALIFEjournal,diStefanoIntention2023},
we read every agent-game against the four canonical rules AllC, AllD,
Tit-for-Tat and Win-Stay-Lose-Shift. The read-out reports which of the four a
ten-round trajectory most \emph{resembles}. It is not an identification of the
strategy a model is running, and no claim in this paper treats it as one.

All four rules are memory-one, so exact matching is a lookup with no learning
in it, and it returns a \emph{set} rather than a label, because over ten rounds
several rules routinely prescribe the same moves. The three provenances, the
agent-games matched exactly by one rule (deduced), by several at once
(ambiguous) and by none (unmatched), are recorded separately and never merged,
and the same matcher yields a continuous quantity needing no labels at all, the
number of rounds violating the nearest rule. Where the rules leave the answer
open a single-layer LSTM~\cite{lstm} arbitrates. Trained on a synthetic corpus alone, which contains no language-model output, it reaches 98.5\% accuracy in distribution and 93.6\% on a held-out
execution-noise condition (appendix~\ref{appendix:classifier}). It ranks only within the set the
rule base returns, so it can never overturn a deduction; it decides alone only
where no rule fits, and its share of the labels is very unevenly spread across
the four rules, which bounds how the composition may be read
(Section \ref{sec:llm_strategies}).

\subsection{Statistical analysis}\label{sec:stats}

The cooperation rate of an agent-game is the share of its ten decisions that
are Option~B; a cell mean averages the agent-games of a cell, and round-level
analyses use the binary decision itself. Because the two agents of a dyad play
each other and the ten rounds inside an agent-game are strongly autocorrelated,
primary inferential intervals were computed by resampling whole dyads with
replacement, 1{,}500 draws, percentile method~\cite{efron1993introduction},
and primary per-model and trend regressions clustered their standard errors on
the dyad. The supplementary scale-by-language robustness model instead clusters
on the documented language-by-game-ID blocks to preserve the matched ten-scale
ladders. Descriptive figure intervals use whole-dyad resampling. Because the range over the ten
scales is a maximum minus a minimum, and so positive by construction, its
bootstrap interval quantifies the precision of the range and not its distance
from zero; the permutation test supplies the test. Because the null is a
theorem, the matching test destroys the association with the payoff scale and
leaves the rest of the design as observed: the scale label is shuffled among
whole dyads, the unit to which the scale was assigned, and the range of the ten
scale means recomputed, 2{,}000 draws, so $0.0005$ is the smallest attainable
$p$ value.

For each model we fitted a logistic regression of the round-level decision on
the payoff scale as a ten-level factor, with prompt language, own persona and
opponent persona as covariates, $\lam = 1$ as reference and standard errors
clustered on the dyad, and tested the nine scale contrasts jointly by a Wald
test (table~\ref{tab:effect}). The variance decomposition is a type-II analysis
of variance on the 250 cell means of the five models reported here, with terms
for model, language and payoff scale and their three two-way interactions and
no three-way term, so the residual on 144 degrees of freedom is the omitted
three-way interaction; the supplementary decomposition is in
appendix~\ref{appendix:variance}. The persona result contrasts the two
sub-unit scales with the remaining eight as a difference of differences with an
interval from the same dyad bootstrap; where a trend across the whole grid is
wanted instead, $\log_{10}\lam$ enters as a continuous predictor. Analyses were
run in Python; the package versions and the seeds are recorded with the
archived pipeline in appendix~\ref{appendix:repro}.


% ===========================================================================
\section{Results}\label{sec:results}

The results below are computed on the 200{,}000 recorded decisions in the
five-model main-text corpus.

\subsection{Cooperation varies under payoff rescaling }\label{sec:payoff_magnitude}

Under normative decision theory and classical game theory, multiplying every
payoff in a normal-form matrix by a positive constant is an axiomatically
neutral transformation. It preserves all preference rankings over lotteries,
best-reply correspondences, and Nash equilibria: an agent playing the game
should not alter its choices whether payoffs are denominated in cents, dollars,
or thousands of pounds.

Yet, across five orders of payoff magnitude, every frontier language model
evaluated violates this scale invariance (see figure~\ref{fig:landscape}a,b,
table~\ref{tab:effect}). The shift in strategic behaviour is substantial: across
the ten-scale ladder, cooperation varies markedly for every model, representing a clear departure from the scale-invariance null (see figure~\ref{fig:landscape}c and
table~\ref{tab:effect}). What classical game theory treats as an inert change of
units produces a fundamental reordering of multi-agent behaviour.

Notably, the models do not share a universal behavioural pattern: each
possesses a distinct machine character that responds idiosyncratically to the
scale of the stakes. Grok exhibits a behavioural strong shift across payoff scales. At micro-stakes, it predominantly  defects; at standard stakes, however, it shifts towards high levels of cooperation, selecting the socially optimal action in the vast majority of interactions. 
GPT climbs sharply
from a hesitant baseline before saturating at a stable cooperative tendency.
Claude’s cooperation rate decreases  at intermediate sub-unit stakes before recovering at larger payoff scales. Gemini maintains a persistently
elevated cooperative tendency throughout, while Qwen drifts steadily in the
opposite direction, defecting progressively more often as the stakes increase.
Because each model follows a distinct behavioural trajectory, machine cooperation cannot be characterised by a single universal behavioural pattern.

The largest departures are concentrated near the low-payoff end of the grid
(table~\ref{tab:effect}). Fractional payoff multipliers often reduce cooperation
compared with intermediate scales, whereas changes from unit stakes to large stakes
have no consistent pattern across models. This result should be interpreted descriptively. The present design does not establish a discrete cognitive
threshold, distinguish numerical magnitude from every feature of notation, or
identify the mechanism behind the response.

\begin{figure}[!htbp]
\centering
\includegraphics[width=\linewidth]{figures/f_landscape.pdf}
\caption{\textbf{Cooperation rates as a function of the payoff scale across five frontier models.}
(a) Model-by-scale cooperation matrix: mean cooperation for each of the five models at each of the ten payoff scales $\lam \in [0.01, 1000]$, with the value shown in each box.
(b) Small multiples per model: the focal model is drawn with its designated colour and marker, with the other four models shown behind in grey for context. Shaded bands mark the sub-unit regime ($\lam \leq 0.1$).
(c) Within-model cooperation range (maximum minus minimum rate across the ten scales) compared against the pooled five-model range (0.154, vertical dashed line). Confidence intervals are 95\% percentile intervals from 1{,}500 whole-dyad bootstrap resamples.}
\label{fig:landscape}
\end{figure}

% Generated by Analysis/scaling/s06_tables.py - do not edit by hand.
\begin{table}[t]
\centering
\caption{\textbf{Cooperation over the ten payoff scales, by model.}
Minimum, maximum and range, with a 95\% confidence interval on the range from a
bootstrap over whole dyads and a permutation test on the payoff-scale label
($0.0005$ is the floor of 2{,}000 draws). The last two columns test the nine
scale contrasts jointly in a per-model logistic regression of the round-level
decision, clustered on the dyad. Each entry pools 200 dyads per scale; ranges
are computed before rounding. Full specification in the electronic
supplementary material.}
\label{tab:effect}
\small
\begin{tabular}{lccccccc}
\toprule
& \multicolumn{2}{c}{cooperation} & & & & \multicolumn{2}{c}{Wald test} \\
\cmidrule(lr){2-3} \cmidrule(lr){7-8}
Model & min & max & range & 95\% CI & $p_{\text{perm}}$ & $\chi^2_9$ & $p$ \\
\midrule
Claude Haiku 4.5 & 0.371 & 0.452 & 0.081 & [0.062,\,0.135] & $0.0025$ & 32.3 & $<0.001$ \\
GPT-5.4 Nano & 0.438 & 0.680 & 0.243 & [0.200,\,0.290] & $0.0005$ & 164.1 & $<0.001$ \\
Gemini 3.5 Flash-Lite & 0.682 & 0.872 & 0.190 & [0.165,\,0.240] & $0.0005$ & 129.0 & $<0.001$ \\
Qwen3 235B-A22B & 0.399 & 0.494 & 0.095 & [0.068,\,0.158] & $0.0005$ & 96.9 & $<0.001$ \\
Grok 4.20 & 0.274 & 0.910 & 0.636 & [0.588,\,0.686] & $0.0005$ & 1227.3 & $<0.001$ \\
\bottomrule
\end{tabular}
\end{table}

\begin{table}[t]
\centering
\caption{\textbf{Persona effect below and above the sub-unit boundary, by
model.} Cooperation under a cooperative persona minus cooperation under a
selfish one, at the two scales where every printed payoff is at most 1
($\lam \leq 0.1$) and at the eight where some payoff exceeds it. Intervals are
95\% confidence intervals from the dyad bootstrap. A sign flip needs both
regime intervals to exclude zero on opposite sides, so a supra-unit interval
straddling zero means the effect is abolished rather than reversed.}
\label{tab:persona}
\footnotesize
\setlength{\tabcolsep}{3pt}
\begin{tabular}{lccccc}
\toprule
& \multicolumn{2}{c}{persona effect, with 95\% CI} & & & \\
\cmidrule(lr){2-3}
Model & $\lam \leq 0.1$ & $\lam \geq 0.25$ & shift & 95\% CI & sign flip \\
\midrule
Claude Haiku 4.5 & $-0.222$ $[-0.258,\,-0.186]$ & $+0.207$ $[+0.189,\,+0.226]$ & 0.429 & [0.389,\,0.471] & yes \\
GPT-5.4 Nano & $-0.135$ $[-0.179,\,-0.091]$ & $+0.018$ $[-0.001,\,+0.037]$ & 0.153 & [0.105,\,0.200] & no \\
Gemini 3.5 Flash-Lite & $-0.456$ $[-0.507,\,-0.406]$ & $+0.224$ $[+0.207,\,+0.242]$ & 0.680 & [0.625,\,0.736] & yes \\
Qwen3 235B-A22B & $-0.911$ $[-0.931,\,-0.893]$ & $-0.679$ $[-0.694,\,-0.664]$ & 0.233 & [0.210,\,0.258] & no \\
Grok 4.20 & $-0.441$ $[-0.485,\,-0.396]$ & $-0.303$ $[-0.326,\,-0.279]$ & 0.138 & [0.090,\,0.185] & no \\
\bottomrule
\end{tabular}
\end{table}

\begin{table}[t]
\centering
\caption{\textbf{Rule agreement and rule distance against the payoff scale,
by model.} Left, the share of agent-games matched exactly by one canonical
rule, below and above the sub-unit boundary, with the coefficient of a logistic
regression of that indicator on $\log_{10}\lam$; right, the mean number of
rounds violating the nearest canonical rule at the smallest and largest scale,
with the same trend by least squares. Both $\beta$ are per decade of $\lam$ and
their signs differ between models. Both cluster on the dyad, and neither block
uses the learned classifier.}
\label{tab:strategy}
\small
\setlength{\tabcolsep}{4pt}
\begin{tabular}{lcccccccc}
\toprule
& \multicolumn{4}{c}{exactly one rule fits (\%)}
& \multicolumn{4}{c}{rounds violating the nearest rule} \\
\cmidrule(lr){2-5} \cmidrule(lr){6-9}
Model & $\lam \leq 0.1$ & $\lam \geq 0.25$ & $\beta$ & $p$
      & $\lam{=}0.01$ & $\lam{=}10^3$ & $\beta$ & $p$ \\
\midrule
Claude Haiku 4.5 & 17.9 & 3.5 & $-0.545$ & $<0.001$ & 1.70 & 2.02 & $+0.073$ & $<0.001$ \\
GPT-5.4 Nano & 9.8 & 7.6 & $-0.038$ & $0.441$ & 2.02 & 2.19 & $+0.038$ & $0.052$ \\
Gemini 3.5 Flash-Lite & 32.2 & 15.6 & $-0.242$ & $<0.001$ & 0.48 & 1.36 & $+0.194$ & $<0.001$ \\
Qwen3 235B-A22B & 72.1 & 53.6 & $-0.267$ & $<0.001$ & 0.27 & 1.14 & $+0.202$ & $<0.001$ \\
Grok 4.20 & 44.6 & 17.1 & $-0.301$ & $<0.001$ & 0.92 & 0.37 & $-0.118$ & $<0.001$ \\
\bottomrule
\end{tabular}
\end{table}


\subsection{Effect sizes vary across prompt languages}\label{sec:language}

Prompt language is not a neutral transmission channel for the game. Both the
baseline level of cooperation and its sensitivity to the stakes diverge
substantially across prompt languages. Figure~\ref{fig:language} shows the
pooled language profiles, while the model-by-language levels and ranges appear
in electronic supplementary table~\ref{tab:S-language}.

We note a caveat related to the comparison of levels rather than the scale effect itself. 
The Arabic and Chinese templates differ from the other prompt stimuli in their
objective wording. They state the four cells as penalties while phrasing the
goal sentence as maximising a reward (Section \ref{sec:protocol}). Because that
wording is fixed within each language, the within-language comparisons still
establish sensitivity to payoff rescaling for the prompt as actually presented.
However, differences in the magnitude or shape of scale sensitivity across
languages cannot be attributed to language alone, because language is partially
confounded with template wording. We therefore interpret the scale-by-language
interaction as a prompt-stimulus interaction, not as evidence for a linguistic
or cultural mechanism.

The central finding is that scale sensitivity itself depends on the prompt
language. The model-specific max-minus-min ranges over the ten scales vary by
up to an order of magnitude within a model (electronic supplementary
table~\ref{tab:S-language}). Figure~\ref{fig:language}b provides the
complementary supra-unit minus sub-unit contrast, and figure~\ref{fig:language}c
shows that language differences remain after removing their mean levels. Across
the main-text corpus, the observed sensitivity ranges from near-invariance to
sweeping behavioural transformations spanning almost the entire probability
scale.

Crucially, these linguistic biases do not point in a uniform direction across
models: each model identifies a different language as its most cooperative
medium and another as its most scale-sensitive. When multilingual evaluations
pool across models, these opposing model-specific effects are attenuated and
can obscure substantial model-by-language heterogeneity. In reality, machine
behaviour depends on the linguistic stimulus used in the prompt. Within this design,
scale sensitivity is a joint property of the model and prompt language rather
than a model-independent invariant. It does not identify a cultural mechanism.

\begin{figure}[!htbp]
\centering
\includegraphics[width=\linewidth]{figures/f_language.pdf}
\caption{\textbf{Cross-lingual scale sensitivity.} (a) Cooperation pooled across the five main-text models for each prompt language and payoff scale; the dashed boundary separates the sub-unit and supra-unit regimes. (b) Language-specific difference in cooperation between supra-unit and sub-unit regimes; intervals show the pooled estimate and small markers the five models separately. (c) Cooperation curves after subtracting each language's mean, showing that differences in scale-response shape remain after level differences are removed.}
\label{fig:language}
\end{figure}

\subsection{Model heterogeneity and the effects of aggregation}\label{sec:heterogeneity}

Averaging across models conceals the real phenomenon through mutual
cancellation. The pooled response curve severely attenuates the true magnitude
of scale sensitivity, shrinking the typical within-model movement substantially
(figure~\ref{fig:landscape}c). Because pairwise correlations between model curves
are predominantly negative (figure~\ref{fig:landscape}b), curves that move in
opposite directions flatten each other out when summed. A pooled curve describes
an artificial average agent that no real model resembles.

Underlying this cancellation is a fundamental divide in policy geometry:
cooperation in these models is not a single Gaussian distribution, but a mixture
over distinct policy regimes (appendix~\ref{appendix:pooling}). Across the corpus,
unconditional boundary play is strongly overrepresented, with extreme atoms
carrying far more weight than expected under uniform play. Crucially, the commitment to boundary
play varies dramatically across models: Claude operates as an interior,
mixed-policy player that rarely locks into extreme  behavioural patterns, whereas Qwen and
Grok function as pure-policy switchers, oscillating between unconditional
defection and unconditional cooperation. Pooling across models conflates how
often a model chooses a pure strategy with which pure strategy it selects.

A formal variance decomposition confirms that model individuality dominates
the game (table~\ref{tab:S-variance}, appendix~\ref{appendix:variance}). Model
identity and the model-by-language interaction account for the majority of the
variance across cells. While the main effect of payoff scale is statistically
detectable, its interaction with model identity carries three times as much
variance, demonstrating that the behavioural impact of stakes is model-specific
rather than shared universally across model endpoints.

In the five-model cell-mean variance decomposition, the shared payoff-scale
main effect is largely carried by Grok. Removing Grok renders that particular
main-effect test non-significant, even though every remaining model continues
to move (appendix~\ref{appendix:variance}). A complementary dyad-level
leave-one-model-out test asks a different question, whether any pooled scale
association remains. Every five-model subset retains a detectable association,
including the subset without Grok (appendix~\ref{appendix:robustness}). The two
analyses are therefore complementary: one tests a shared main effect after
accounting for heterogeneity, whereas the other tests any residual pooled
association.

\subsection{Payoff scale modulates persona-induced cooperative behaviour}\label{sec:persona}

System prompts designed to steer agent behaviour do not operate independently of
game stakes. The behavioural difference between cooperative and selfish persona
instructions changes dramatically across the payoff ladder (figure~\ref{fig:persona},
table~\ref{tab:persona}). This modulation differs between the two sub-unit
scales and the eight larger scales, where shown payoffs include values above
one (Section \ref{sec:gameImportance}). Across all five models, the persona effect
changes across these two regimes.

In multiple frontier models, the persona instruction reverse their behavioural effect  below the
unit threshold. In Gemini and Claude, the persona effect reverses sign: under fractional stakes, both models exhibit higher cooperation when explicitly instructed to behave selfishly than when instructed to behave cooperatively. In GPT, the persona
effect is  extinguished at fractional scales rather than reversed,
while the remaining models maintain their pattern. This result is a
boundary-aligned descriptive effect. It does not identify why the prompt
instruction changes direction at these scales.

The main grid does not separate numerical magnitude from all aspects of numeral
formatting. We therefore treat this regime split as descriptive and do not
attribute it to a specific representational mechanism.

A separate class of models inverts the persona instruction across the entire
ladder (table~\ref{tab:persona}). Qwen and Grok cooperate substantially more
under the selfish persona than under the cooperative persona at every evaluated
scale. These models show an opposite response to the persona wording. The data
do not establish why they do so. This result underscores that role prompts
cannot be assumed to steer multi-agent interactions in the intended direction
without empirical verification across payoff contexts
(appendix~\ref{appendix:persona_full}).

\begin{figure}[!htbp]
\centering
\includegraphics[width=\linewidth]{figures/f_persona.pdf}
\caption{\textbf{Persona effect and sub-unit inversion against the payoff scale.}
The persona effect is the cooperation rate under the cooperative persona minus that under the selfish persona.
(a) Difference in cooperation rate across scales $\lam$, with the sub-unit regime ($\lam \le 0.1$) shaded where all five models cooperate more under the selfish prompt.
(b) Regime means comparing sub-unit ($\lam \le 0.1$, open circles) with supra-unit ($\lam \ge 0.25$, filled circles) regimes (table~\ref{tab:persona}).}
\label{fig:persona}
\end{figure}

\subsection{ Payoff scale alters strategic composition, not merely cooperation rates}\label{sec:llm_strategies}

Two models displaying identical cooperation rates can be executing fundamentally
different behavioural strategies/rules. To understand the behavioural rules underlying raw
cooperation rates, we read every agent trajectory against the four canonical
memory-one strategies of evolutionary game theory (Section \ref{sec:readout}). What
this read-out captures is behavioural resemblance rather than exact algorithmic
identity: across the corpus, the majority of games match no textbook rule
cleanly (table~\ref{tab:strategy}, appendix~\ref{appendix:readout_full}).

As payoff stakes increases, play moves farther from the four canonical rules in
most frontier models. Across the scale ladder, the proportion of games adhering
strictly to a single canonical rule declines substantially, while unmatched
play accumulates (table~\ref{tab:strategy}). This loss of canonical-rule
conformity is pronounced in Claude, Grok, Qwen, and Gemini, while GPT maintains
steady adherence. The read-out does not establish what rule generates unmatched
play; it only shows that it is farther from the four canonical rules.

A label-free continuous metric confirms this behavioural drift: the number of
rounds departing from the nearest canonical rule increases significantly with
scale across most models (table~\ref{tab:strategy}). As stakes multiply,
trajectories depart more often from the textbook templates. Only Grok reverses
this trend, moving towards extreme pure-policy play at high stakes.

The qualitative composition of strategies also shifts dynamically, driving
divergent realignments across models (figure~\ref{fig:mix}). In GPT, expanding
stakes prompts a steady migration from unconditional defection toward
unconditional cooperation. In Gemini, the opposite shift occurs: unconditional
cooperation yields ground to more Tit-for-Tat-like and Win-Stay-Lose-Shift-like
attributions. Grok undergoes an outright
reversal, swapping rampant defection at micro-stakes for overwhelming cooperation
at high stakes.

Critically, these strategy shifts reflect qualitative resemblances rather than
exact implementations of them. Provenance auditing demonstrates that while unconditional
cooperation and defection are often matched exactly, nearly all attributions of
Tit-for-Tat and Win-Stay-Lose-Shift represent best-fit approximations to noisy,
unmatched play (figure~\ref{fig:mix}c, appendix~\ref{appendix:readout_full}).
When models appear to increase their Tit-for-Tat share, they are not adopting the
classic reciprocal algorithm; they are generating complex, imperfect trajectories
whose overall statistical contours happen to lean in that direction.

\begin{figure}[!htbp]
\centering
\includegraphics[width=\linewidth]{figures/f_strategy_mix.pdf}
\caption{\textbf{Strategy attribution across payoff scales, composition shift, and provenance breakdown.}
(a) Stacked area share of canonical memory-one strategy attributions (AllC, AllD, TFT, WSLS) across the ten payoff scales for each of the five models.
(b) Net strategy share changes between extreme scales ($\lam = 10^{3}$ minus $\lam = 0.01$) across the four canonical rules with total variation distance $D_{\mathrm{TV}}$ in parentheses.
(c) Strategy read-out provenance: exactly deduced (single canonical match), ambiguous (multiple matching rules), and unmatched (LSTM attribution on trajectories consistent with no standard rule). Over the corpus, $94.6\%$ of Tit-for-Tat and $97.1\%$ of Win-Stay-Lose-Shift labels are attributions on unmatched play, establishing that labels indicate resemblance rather than faithful execution.}
\label{fig:mix}
\end{figure}

\subsection{Where payoff scale  acts: the evolutionary baseline}\label{sec:egt_pointer}

While payoff scale is theoretically inert in normal-form games, evolutionary
game theory provides a foundational framework where stakes legitimately
matter, offering a rigorous theoretical benchmark for language model behaviour.

In deterministic replicator dynamics, scale is completely neutral. Multiplying
the payoff matrix by a positive constant simply scales the velocity of the
vector field: orbits follow identical geometric trajectories and all rest points
remain invariant (appendix~\ref{appendix:egt_simplex}, figure~\ref{fig:app_simplex}).

In finite populations with stochastic imitation \citep{nowak2004emergence}, however, payoff scale acts
directly through the selection intensity. Under pairwise Fermi imitation, scaling payoffs
amplifies the competitive advantage of defection over cooperation
(appendix~\ref{appendix:egt}). At micro-stakes, selection is weak and all four
canonical rules coexist near genetic drift. As stakes expand, selection pressure
intensifies: unconditional defectors strongly invade cooperative strategies until
AllD achieves near-total fixation (figure~\ref{fig:egtllm}a). Combining
evolutionary dynamics with language-model evaluations provides a principled
lens for multi-agent governance~\cite{balabanova2025media,buscemi2025llms,xie2026evolutionary,han2026social}.

At the pooled five-model level, the AllD label share moves in the opposite
direction to the finite-population evolutionary baseline as $\lam$ increases
(figure~\ref{fig:egtllm}). This pooled reversal is driven primarily by GPT and
Grok; per-model label trajectories remain heterogeneous. As payoffs grow into
standard integer ranges, the pooled language-model corpus becomes more
cooperative, whereas the evolutionary reference becomes more dominated by
defection (figure~\ref{fig:egtllm}b,c).

This direction is stable across the eighty evolutionary parameter settings
reported in the appendix: the stationary AllD share is higher at the largest
than at the smallest scale in every setting, and the rise is strictly monotonic
in 61 of 80 settings. None reproduces the direction of the pooled empirical
trajectory.
We therefore treat the comparison as a descriptive reference gap rather than as
evidence for a particular mechanism. Evolutionary dynamics describe selection
among rules, whereas the corpus describes fixed language-model agents.

The gap is a descriptive difference between selection among rules and fixed
language-model agents. The language conditions also vary the stimulus, so this
comparison does not identify a cultural or cognitive mechanism.

\begin{figure}[!htbp]
\centering
\includegraphics[width=\linewidth]{figures/f_egt_vs_llm.pdf}
\caption{\textbf{The evolutionary baseline and the corpus move in opposite
directions.} (a) The baseline's stationary distribution over the four canonical
rules across the ten payoff scales, drawn as a stacked band. (b) The corpus's
strategy mix on the same axis, drawn identically. (c) The AllD share alone,
both accounts on one axis, with the gap shaded. Baseline at $Z = 100$,
$\beta = 0.1$, $\epsilon = 0.05$, ten rounds; the corpus pooled over five
models, five languages and four persona pairings.}
\label{fig:egtllm}
\end{figure}

\subsection{Agents condition more strongly on their own past actions}\label{sec:conditioning}

The evolutionary model rests on an explicit behavioural mechanism that can be
evaluated directly in the data. In the classical account, cooperation survives
invasion by defectors through \emph{reciprocity}: conditioning one's choice on
what the partner just did, retaliating against defection and rewarding
cooperation. Tit-for-tat is the prototypical exemplar of this logic.

Because the previous round consists of an action pair (own action, partner
action), the four conditional cooperation probabilities decompose cleanly into
a base level and three orthogonal contrasts. \emph{Reciprocity} represents the
partner main effect: the difference in cooperation after a partner cooperates
versus after a partner defects. \emph{Persistence} represents the own main
effect: the tendency to repeat one's own previous action. \emph{Matching}
represents the interaction effect characteristic of Win-Stay-Lose-Shift.
Textbook strategies occupy precise locations in this geometric space: Tit-for-Tat
lies entirely on the reciprocity axis, Win-Stay-Lose-Shift on the matching axis,
and unconditional rules at the origin (figure~\ref{fig:strategyspace}a).

Decomposing agent play reveals that own-action conditioning exceeds partner
conditioning at the model level (figure~\ref{fig:strategyspace}). Across all
five models, the persistence contrast is larger than the reciprocity contrast,
although the margin is modest for Claude and substantially larger for the other
models. Reciprocity remains positive in four models and is negative in Qwen. An
unsupervised principal component analysis on raw transition
vectors independently recovers this identical geometry without any prior
knowledge of the contrasts (figure~\ref{fig:strategyspace}c).

The contrast estimates are precise at the dyad level. Whole-dyad bootstrap
intervals for reciprocity are positive for Claude, GPT, Gemini~3.5, and Grok,
and negative for Qwen; the corresponding persistence intervals are positive for
all five models (electronic supplementary table~\ref{tab:conditioning-ci}). Thus the
headline comparison is one of relative strength, not a claim that partner
information is absent.

This finding narrows one candidate explanation for the divergence from the
evolutionary baseline. In these data, agents condition more strongly on their
own preceding moves than on their partner's behaviour. The contrast does not
show that partner information is absent, nor does it identify the causal basis
of cooperation.

A technical caution governs this decomposition. The naive difference between
disagreement states, $p_{CD} - p_{DC}$, is frequently mistaken for a reciprocity
metric in the literature. In reality, it equals persistence minus reciprocity
identically, confounding the two main effects and falsely indicating
anti-reciprocal behaviour where the true partner effect is positive.

\begin{figure}[!htbp]
\centering
\includegraphics[width=\linewidth]{figures/f_strategy_space.pdf}
\caption{\textbf{Own-action conditioning exceeds partner-action conditioning at the model level.}
(a) Cells in the reciprocity by persistence plane, with the canonical rules at
their exact coordinates; AllC, AllD and win-stay lose-shift all have both main
effects zero and so coincide at the origin. (b) The three contrasts per model,
against the value each canonical rule would give, with 95\% percentile intervals
from whole-dyad bootstrap resampling. (c) The first two principal components of
the raw five-vector, which is given no knowledge of the contrasts and recovers
the same geometry. Panels (a) and (c) use model-by-scale-by-language cells; a
predecessor observed fewer than thirty times yields no estimate and the cell is
dropped, leaving 212 of 250.}
\label{fig:strategyspace}
\end{figure}


\section{Concluding remarks and discussion}\label{sec:discussion}

Positive scaling of payoffs is theoretically behaviour-preserving in the (repeated) Prisoner's Dilemma: they leave the strategic structure of the game unchanged and therefore imply no change in predicted outcomes. Against this, our results show systematic departures from scale invariance across all five frontier language models examined. These departures are not reducible to a homogeneous shift in baseline cooperation. Their magnitude varies across prompt languages, their direction and shape differ across models, and they are substantially attenuated by aggregation. 
The resulting pooled trajectory is therefore not an adequate representation of any individual model's behavioural response. 
Relative to an evolutionary game baseline computed on the same payoff matrices, the pooled corpus also exhibits a declining share of AllD-like behaviour as payoffs increase, moving in the opposite direction from the baseline prediction.
This comparison quantifies the distance between the incentives encoded by the game and the behaviour produced by the agents, without by itself identifying the mechanism responsible for that distance.

If scale invariance fails, a reported cooperation rate can no longer be interpreted as a scale-independent property of the model. This implication is central to machine behaviour, which argues that autonomous systems should be studied empirically, using the instruments of the behavioural sciences, because their behaviour cannot be inferred from their construction alone~\cite{rahwan2019machine}. Yet behavioural experiments face a familiar difficulty: most manipulations do not imply a precise predicted effect, so any observed change must be interpreted relative to a comparison condition chosen by the experimenter. The invariance classes of decision theory provide an unusual exception. They specify manipulations for which the predicted behavioural effect is exactly zero. A violation of such an invariance is therefore especially informative: it does not merely show that behaviour changes under a different experimental condition, but that it changes under a transformation that should be behaviourally inert.

This point has consequences for several strands of machine-behaviour research touched by this study. Persona prompting is the clearest case in which the experimenter addresses the system in terms of its social role; our results show that whether such an instruction takes hold can depend on the numerical units in which the payoff matrix is presented. The implication is equally direct for models of mixed human--machine populations, which treat the cooperativeness of artificial agents as a parameter and ask where the population settles~\cite{santos2018social,han2011intention,han2013good,perc2017statistical,paiva2018engineering}. A cooperation rate partly determined by the experimenter's choice of units is not the scale-independent quantity such models require.


The methodological implication of our findings is significant. Cooperation rates elicited from language models are often treated as properties of the model under evaluation, conditional on a prompt and a game form. The present evidence indicates that they are instead joint properties of the model, the prompt language, the persona instruction, and the numerical scale at which the payoff matrix is expressed. The last of these is especially important because it is usually treated as theoretically irrelevant. 
In standard game-theoretic analysis, multiplying every payoff-matrix entry by a positive constant changes units but not incentives. 
For the models studied here, however, this formally irrelevant choice systematically changes observed behaviour. Thus, the payoff scale cannot be regarded as a neutral implementation detail in empirical studies of machine cooperation.

This conclusion is strengthened by the structure of the heterogeneity. The observed scale effects are not merely noisy fluctuations around an invariant response. They differ by model, interact with prompt language, moderate the behavioural effect of persona instructions, and alter the composition of inferred strategic patterns. These features are difficult to reconcile with sampling variability alone. They suggest that the models are sensitive to aspects of the numerical representation of the game that are not part of its strategic equivalence class. However, because the design manipulates payoff scale rather than the internal representation of number, notation, or stakes, the results should be interpreted as evidence of invariance violation rather than as evidence for any particular psychological or computational mechanism.

%The aggregation results are especially consequential. In the pooled data, a statistically significant main effect of payoff scale remains, but it explains substantially less variation than the model-by-scale interaction. Moreover, the pooled range understates the typical model-specific range by a factor of 1.62. The pooled mean therefore summarizes a behavioural trajectory that no individual model need follow. This is not simply a loss of precision that could be remedied by increasing the number of dyads. When model-specific scale responses differ in shape or direction, additional observations tighten confidence intervals around an aggregate that is itself substantively ambiguous. The appropriate unit of analysis is therefore the model-specific response surface, with prompt language and persona treated as factors that may themselves condition scale sensitivity.

%This point bears directly on benchmarking. A single cooperation rate measured at one payoff scale cannot be assumed to generalize across equivalent numerical representations of the same game. Two studies using the same model, prompt template, and game form could obtain markedly different cooperation rates if their payoff matrices differ only by a constant multiplier. Conversely, pooled comparisons across models may obscure behaviourally meaningful scale effects when models respond in opposing directions. Benchmark reports should therefore record the numerical scale of the payoff matrix and, where possible, include a scale-sensitivity analysis alongside any headline cooperation rate. Such analyses are particularly important when cooperation rates are subsequently used as inputs to population models, evolutionary simulations, or safety evaluations.

We also qualify the interpretation of persona effects. Persona instructions are often treated as relatively direct interventions on behavioural disposition, such as inducing more cooperative or more selfish play. The present results show that their effect may depend on the payoff scale at which the interaction is framed. In several LLMs, persona instructions reverse their behavioural effect below the unit threshold: under fractional stakes, models instructed to behave selfishly may cooperate more than those instructed to behave cooperatively. This pattern does not imply that the models possess stable preferences that are being inverted by small payoffs. Rather, it shows that the behavioural consequence of a verbal instruction is conditional on the numerical context in which that instruction is applied. Persona effects are therefore not separable from payoff representation.

The strategy-composition analysis extends the same conclusion beyond marginal cooperation rates. Payoff scale changes not only how often cooperation occurs, but also the kinds of action sequences that the models generate. These shifts indicate qualitative resemblances to canonical repeated-game strategies rather than exact implementations of them. The distinction is important. A ten-round trajectory cannot establish that a model is executing Tit-for-Tat, Win-Stay-Lose-Shift, Grim Trigger, or any other formal strategy. What can be inferred is that some trajectories are more consistent with the behavioural signatures of such strategies than others, and that the prevalence of these signatures changes with payoff scale. Scale sensitivity therefore operates at the level of sequential structure as well as aggregate cooperation.

%The relationship between these findings and existing work on language-mediated cooperation is complementary. 

Prior studies have examined how cooperation changes under moral framing, pre-play communication, institutional mediation, or other contextual interventions \cite{buscemi2025strategic,lu2024llms,mao2025alympics,xie2026evolutionary}. Those manipulations can  alter the game as interpreted by the agent, because they change the semantic environment in which the decision is made. 
Payoff rescaling is different: under the standard formal model, it preserves the game. The fact that behaviour nevertheless changes indicates that language models may not respect the invariance classes assumed by game-theoretic representations. More broadly, if model behaviour is governed by a utility-like object that combines numerical payoffs with linguistic or representational features, then scale invariance will generally fail unless the numerical and non-numerical components transform together. Any theoretical account of machine play in strategic settings must therefore specify which transformations it regards as behaviourally equivalent.

Several possible mechanisms could produce the observed departures. Larger numerical magnitudes may be interpreted as higher stakes, even though positive rescaling does not change strategic incentives. Different numeral strings may also have different tokenization profiles, different frequencies in training data, and different learned associations. 
Fractional payoffs may be especially unusual in contexts framed as rewards, causing models to treat them as qualitatively distinct rather than merely smaller quantities. Alternatively, the effect may arise from interactions among magnitude, notation, language, and instruction following. 
The present design cannot distinguish among these explanations. It does not independently manipulate tokenization, hold magnitude fixed while varying notation across a broad set of formats, or vary the semantic framing of stakes. Accordingly, the results identify a robust behavioural phenomenon but not its causal source.

Several limitations follow from this design. First, the study examines five models, one repeated game, one family of prompt templates, and a finite set of ten payoff scales. The results therefore support claims about the presence and structure of scale sensitivity in this corpus, not universal claims about all language models or all strategic environments. Second, the language comparison is bounded by the prompt templates used. Two templates describe payoffs as penalties while retaining a reward-maximization objective. Because this wording is fixed within each language across payoff scales, it does not undermine the within-language rescaling comparison. It may, however, interact with numerical magnitude and therefore limits any attribution of cross-language differences to language alone.

% Third, the statistical reference points used here quantify departures relative to the variability observed in the corpus. The permutation tests and dyad-level bootstrap establish that the observed ranges are unlikely under the relevant exchangeability assumptions, but they do not provide a full hierarchical model of replicate, dyad, prompt, and game-level dependence. Clustering at the dyad level addresses important dependence among actions within repeated interactions, yet a more complete design would include independent replications and mixed-effects models with random intercepts or slopes for replicate and game. The independent repeat of Gemini 3.5 at four scales, reported in the appendix, provides useful evidence on repeatability for that arm, but it should not be interpreted as a general test-retest estimate for the full corpus.

Third, the strategy analysis is necessarily approximate. Strategies are latent decision rules, whereas the data consist of finite action sequences. Ten rounds provide limited information for distinguishing among memory-one rules, especially when execution noise is present. The classifier was validated up to 20\% execution noise, corresponding to two deviations in a ten-round sequence, and some models operate near that boundary on average. The resulting labels should therefore be read as descriptive summaries of trajectory resemblance, not as evidence that the models implement the corresponding strategies. Longer horizons, richer strategy classes, and mixture-based estimators would be needed to make stronger claims about strategic implementation.
Finally, the evolutionary baseline should be interpreted as a theoretical reference distribution rather than as a fitted model of language-model behaviour. It maps stochastic behaviour onto a symmetric per-action error process, which is analytically useful but necessarily coarse. Prompt sensitivity, model heterogeneity, language effects, end-game reasoning, and scale-dependent instruction following are not well captured by symmetric action flips. The baseline is therefore valuable precisely because it clarifies the direction and magnitude of the discrepancy between formal incentive predictions and observed play, not because it offers a complete generative account of the data. 

Overall, our study shows that payoff scale, although formally irrelevant under standard game-theoretic invariance, systematically affects machine cooperation. The effect is model-specific, language-dependent, and consequential for both cooperation rates and inferred strategic composition. These findings call for caution in interpreting single-scale cooperation benchmarks and suggest that scale sensitivity should be reported as a standard robustness measure in studies of language models in games. More generally, they indicate that empirical work on machine strategic behaviour cannot assume that formal equivalences in the game are automatically respected by the agents that are prompted to play it.


% ===========================================================================
% \section{Discussion}\label{sec:discussion}

% \subsection{Payoff-scale invariance is violated}\label{sec:disc_invariance}

% The measurement reported here is unusual in having a null that is a theorem.
% That property lets a movement in cooperation be read as a departure from an
% invariance rather than as the difference between two conditions whose expected
% gap was never specified. The predicted difference between a matrix and the same
% matrix with larger cells is not small, or contested. It is zero.

% All five models depart from it, and none is close to invariant. What makes the
% departures hard to set aside is not their size but their structure: their size
% depends on the language the game is posed in, they gate the persona instruction
% at a boundary of magnitude, and they change what the transcripts look like as
% sequences. Sampling variability alone is difficult to reconcile with a pattern
% of that shape. The models appear to be responding to something in the matrix
% other than the game it encodes, though this design does not identify what.

% \subsection{Heterogeneity, and why pooled designs miss it}\label{sec:disc_pooling}

% The pooling result deserves emphasis, because it is the mechanism by which an
% effect of this size could have gone unnoticed. The cause is arithmetic rather
% than statistical: the shapes disagree, and several pairs are anticorrelated.
% Heterogeneous effects that cancel are invisible to pooled designs by
% construction.

% In this panel the cancellation is partial rather than total. A main effect of
% the payoff scale does survive the averaging, and it is significant, but it
% carries a third of the variance its interaction with the model carries, and the
% pooled range understates the typical model by a factor of 1.62: a statement
% about an average model that nobody deploys. A conclusion that changes when one
% model joins or leaves a panel is not a conclusion about language models. In the
% cell-mean shared-main-effect test, removing Grok makes the scale term
% non-significant; the complementary dyad-level leave-one-model-out association
% remains detectable. We therefore do not defend the pooled average as a single
% model-level quantity.

% The pooled estimate is therefore not a weaker version of the per-model result
% but a different quantity, and it is not a question of power either: more dyads
% or more models only tighten the interval around an average of curves that point
% in different directions. The remedy is a per-model analysis, not a larger
% sample. The same logic runs one level deeper, since the effect depends on the
% prompt language as well as on the model, so a per-model number pooled over
% languages is the same trick played on a smaller stage. The heterogeneity is not
% a nuisance to be averaged away. It is the finding.

% \subsection{Implications for benchmarking and governance}\label{sec:disc_benchmarking}

% A cooperation rate published for a language model is currently understood as a
% property of the model. On this evidence it is a joint property of the model,
% the units the matrix was written in, and the language the game was posed in.
% The units are the awkward member of that list, because they are the one nobody
% is asked to record: framings get described, personas get quoted, prompts get
% reproduced in appendices, while the payoff scale is fixed once, silently, on
% the authority of a theorem saying the choice is empty. Two laboratories running
% the same protocol on the same model, with matrices differing only by a constant
% rescaling factor, could report radically conflicting cooperation rates spanning
% more than half the probability scale.

% We would suggest that a scale-sensitivity measurement accompany a published
% cooperation rate, in the way an error bar does, and be measured per model, most
% of all where such rates leave the benchmark and become the parameters of a
% population model. The same argument applies to safety auditing. An audit that
% runs in English at one fixed payoff matrix tests a single point on a much
% larger behavioural surface, and the two cheapest stress dimensions to add are
% the two measured here: a sweep over the payoff scale, and a parallel run in the
% prompt's non-English variants. Our central claim is complementary to a
% parallel line of work asking how external \emph{mechanisms} change cooperation,
% through contracts and mediation~\cite{tewolde2026coopeval}, moral
% framing~\cite{backmann2025moralsim} or pre-play
% communication~\cite{loreheydari2026communication}. Payoff scale joins that list
% as an intervention on the numerical structure of the game rather than on the
% institutions around it. The behavioural profiles it
% produces are also the input that models of hybrid human-machine populations
% consume~\cite{han2026social}, and a profile that depends on the experimenter's
% units is not the quantity such a model needs.

% \subsection{Interpretation and possible mechanisms}\label{sec:disc_mechanisms}

% Two observations constrain what an explanation must account for. The boundary
% gating the persona instruction appears to be one of magnitude rather than of
% notation, a separation only as strong as the single scale that supports it
% (Section \ref{sec:persona}); and the payoff scale moves how closely play tracks any
% canonical memory-one rule at all, and shifts the mix of labels with it. The
% scale is not turning a single dial marked cooperativeness; it is changing the
% shape of the trajectory.

% Several accounts are consistent with both. The shown magnitudes may be read
% as stakes, which is exactly what a rescaling does not change; numerals of
% different magnitudes are tokenised differently and are met at very different
% frequencies in training text, so a representation that changes with the digit
% string could change what follows from it with no notion of stakes
% involved~\cite{razeghi2022impact}; training text may associate large quantities
% with caution; and sub-unit payoffs may fall outside the range in which
% payoff-like numbers are ordinarily met. We can order none of these. The design
% measures departures from an invariance; it does not manipulate tokenisation,
% vary notation at fixed magnitude beyond the one contrast above, or move the
% framing that supplies the stakes. Nothing here should be read as identifying a
% mechanism.

% There is a constructive reading as well as a cautionary one. A programme is
% under way to widen the utility function beyond the economic consequences of an
% action, so that the language describing the action enters it directly, with the
% replicator dynamics named as the destination~\cite{capraro2024language}. Our
% result sits one level below that programme rather than beside it. The evidence
% for language-based preferences comes from manipulations that could legitimately
% have mattered, since a label carries connotation, whereas a positive rescaling
% carries none and is inert by theorem, and the behaviour still moves. A utility
% that mixes a payoff with a score computed from language is no longer scale-free
% in the payoff argument unless the two terms are rescaled together, so before
% such an object is carried into an evolutionary dynamic its invariance class has
% to be stated. What these data say is that at least one class of agents does not
% behave as though the payoff argument were scale-free, and that anything built
% to describe them will have to account for it.

% \subsection{Limitations and future work}\label{sec:disc_limitations}

% This design does not identify the mechanism inside the model that produces the
% departures it measures across a broad but finite grid. Five models cannot
% support claims about model families or about parameter count, and what
% generalises here is that the shapes differ rather than what any one of them is.
% The corpus is one game, one prompt template family and one horizon, so whether
% the same curves appear in other dilemmas or in longer play is untested.

% One feature of the stimuli bounds the language comparison. Two of the five
% templates state the payoffs as penalties while phrasing the goal as maximising
% a reward, and we reproduced rather than repaired them so that the corpus
% matches the published protocol. The wording is fixed within each language
% across the ten scales, so it does not invalidate the within-language rescaling
% comparison. It may nevertheless interact with numerical magnitude; consequently,
% cross-language differences in scale sensitivity cannot be attributed to language
% alone.

% Our reference for calling a departure large is the sampling variability of
% these data, a bootstrap over dyads and a permutation null. A targeted
% independent repeat of Gemini 3.5 at four scales is reported in the appendix;
% it is useful evidence about repeatability in that arm but is not a test-retest
% floor for the main-text corpus.
% Clustering on the dyad addresses the dependence between the two agents of a
% game and between the rounds within an agent-game, but it does not turn the
% design into a mixed-effects model with random intercepts for replicate and
% game, which would be the more complete treatment of a nested design.

% The strategy read-out carries a further caveat, which is why its provenance is
% reported wherever a label appears (Section \ref{sec:readout}). A strategy is a
% decision rule and not a decision, so it is never observed directly, and even
% where the resident rule is fixed by construction, separating the memory-one
% rules from observed play has been shown to need a long run of
% observations~\cite{kim2021win}. Our trajectories are ten rounds long. The claim
% that play becomes less rule-governed rests on exact matching and no learned
% component. The classifier was validated to $20\%$ execution noise, which is two
% rounds in ten, and in two of the five models play sits about that far from the
% nearest canonical rule on average, so the read-out is working at the edge of
% the range it was validated on. The compositional claim uses labels the
% classifier supplies wherever the rules leave the answer open, is unevenly
% anchored across the four rules, and is the weaker of the two. Any reading of it
% as evidence that these models implement Tit-for-Tat would be unsupported by our
% data. Richer strategy vocabularies, including zero-determinant and extortionate
% play~\cite{pressDyson2012,hilbe2013memory} and mixture estimators over longer
% horizons~\cite{fontana2025nicer,inferenceStrategies}, are the natural extension.

% Finally, the evolutionary baseline maps the stochasticity of a language model
% onto a symmetric per-action error rate, which is a clean theoretical object but
% a coarse abstraction: prompt sensitivity, end-game reasoning and the
% heterogeneity documented above are not well captured by symmetric flips. The
% baseline should be read as a reference distribution rather than as a fitted
% model, and the noise calibration reported in
% appendix~\ref{appendix:egt_calibration} is a first step rather than an
% estimate.


% % ===========================================================================
% \section{Conclusion}\label{sec:conclusion}

% A positive rescaling of a payoff matrix is the rare manipulation whose
% predicted effect is not merely small but exactly zero. Five frontier language
% models were played across ten such rescalings, and every one of them moved. The
% movement is not a uniform bias a calibration could absorb: its shape differs by
% model and its size by prompt language, so most of it disappears from any
% analysis that averages models, and what survives the averaging depends on which
% models happen to be in the average. That is how a departure of this size can
% sit inside a literature that has not reported it. Read against an evolutionary
% baseline computed on the same matrices, the pooled corpus's AllD share moves in
% the opposite direction as the payoffs grow. This measures rather than explains
% the distance between what the game rewards and what these agents do.

% What follows is a question about design rather than a verdict on any model. A
% cooperation rate measured on a language model is reported as though the matrix
% behind it were an arbitrary and therefore harmless choice. It is arbitrary, and
% it is not harmless. The same argument extends to every invariance the theory
% supplies, from shifts of the payoffs to relabelling of the actions, and each is
% an equally cheap test of whether an agent is playing the game it was shown.


%%%%%%%%%% End matter %%%%%%%%%%

\section*{Ethics statement}
This work did not require ethical approval from a human subject or animal welfare committee.

\section*{Data accessibility}
The game transcripts, prompt templates, trained strategy classifier and analysis code supporting this study are available at
\url{https://github.com/trungkiet2005/prisoner-dilemma-llm}. A versioned archive and DOI will be added on publication.

\section*{Declaration of AI use}
Large language models are the object of study in this article. Generative AI was used only as an assistant to identify and fix bugs in code and to edit English grammar.

\section*{Author contributions}
T.A.H.: conceptualisation, methodology, supervision, project administration, writing (original draft), writing (review and editing).
T.-K.H.: conceptualisation, software, formal analysis, visualisation, writing (original draft).
D.-S.D.-M.: software, data curation, formal analysis, writing (original draft).
T.-B.C., P.-H.L., H.-D.N., P.-Q.N.-L., M.-L.N.-V., H.-P.P., P.-H.P., T.-K.T., C.-N.T., H.T. and G.-T.T.-L.: investigation, data curation, validation, writing (review and editing).
A.B.: methodology, resources, writing (review and editing).
L.H.T.: conceptualisation, methodology, supervision,  writing (review and editing).
T.A.H., T.-K.H. and D.-S.D.-M. contributed equally to this work.
All authors gave final approval for publication and agree to be held accountable for the work performed therein.

\section*{Competing interests}
We declare we have no competing interests.

\section*{Funding}
T.A.H. is supported by EPSRC (grant EP/Y00857X/1) and also acknowledges travel support from HCMUT\text{-}VNUHCM through the Adjunct Professorship scheme.

%%%%%%%%%% Bibliography %%%%%%%%%%
\bibliographystyle{unsrtnat}
\bibliography{mybib}

\end{document}
